\documentclass{article}

% Evaluations & Datasets track — most appropriate for a benchmark paper
% Switch to \usepackage[eandd, final]{neurips_2026} upon acceptance
\usepackage[eandd]{neurips_2026}

\usepackage[utf8]{inputenc}
\usepackage[T1]{fontenc}
\usepackage{hyperref}
\usepackage{url}
\usepackage{booktabs}
\usepackage{amsfonts}
\usepackage{amsmath}
\usepackage{nicefrac}
\usepackage{microtype}
\usepackage{xcolor}
\usepackage{graphicx}

% ---------------------------------------------------------------
% TODO: Replace with your actual title
\title{Probing the Failure Boundary of Metacognitive Calibration in Large Language Models}

\author{%
  % TODO: Fill in when de-anonymizing
  % Your Name \\
  % IIIT Bangalore \\
  % \texttt{your@email.com}
}
% ---------------------------------------------------------------

\begin{document}

\maketitle

% ---------------------------------------------------------------
\begin{abstract}
	% TARGET: ~150 words. Write this LAST.
	%
	% Structure:
	%   [Gap]      Current benchmarks measure average calibration but do not
	%              systematically probe where and why calibration breaks.
	%   [Approach] We introduce [NAME], a benchmark that stress-tests metacognitive
	%              ability at the failure boundary — tasks engineered to sit at the
	%              edge where model confidence diverges from epistemic warrant.
	%   [Design]   Tasks span three regimes: complete information, withheld
	%              information, and deliberately misleading information.
	%   [Finding]  Experiments across N frontier models reveal that [key finding].
	%   [Claim]    Our benchmark exposes failure modes invisible to existing evals.

	\textbf{TODO: Write abstract last.}
\end{abstract}


% ---------------------------------------------------------------
\section{Introduction}
% TARGET: ~400 words. Write this SECOND (after Section 3).
%
% Paragraph 1 — The big picture problem
%   Why does metacognition matter for AI systems?
%   What goes wrong when a model is miscalibrated?
%   (overconfident hallucination, silent failure on hard tasks)
%
% Paragraph 2 — What current benchmarks do
%   Name 2-3 existing benchmarks (e.g. Kadavath 2022, Xiong 2024,
%   AAAI 2025 decoupling framework). What they measure. Where they stop.
%
% Paragraph 3 — The specific gap
%   Average calibration \neq calibration at the boundary.
%   Models can look well-calibrated on aggregate while being systematically
%   wrong in the regime that matters most: ambiguous, OOD, misleading inputs.
%
% Paragraph 4 — Our contribution (bullet list or prose)
%   - We introduce [BENCHMARK NAME]
%   - Three task regimes (defined in Section 3)
%   - Novel elicitation protocol (Section 3.3)
%   - Regime-stratified ECE scoring (Section 4)
%   - Released as Kaggle Community Benchmark
%
% Paragraph 5 — Paper roadmap (1-2 sentences)

\textbf{TODO: Write introduction after Section~\ref{sec:design} is solid.}


% ---------------------------------------------------------------
\section{Background and Related Work}
% TARGET: ~300 words

\subsection{Calibration in Machine Learning}
% Guo et al. 2017 — ECE definition, temperature scaling
% Brier score as a proper scoring rule
% One tight paragraph.

\textbf{TODO: }

\subsection{Metacognition in Large Language Models}
% What has been measured so far:
%   Kadavath et al. 2022 — "Know What You Don't Know"
%   Xiong et al. 2024   — "Can LLMs Express Uncertainty?"
%   AAAI 2025           — Decoupling metacognition from cognition
%   OpenReview 2025     — Functional metacognitive abilities (difficulty
%                         awareness, confidence adjustment, strategy revision)
%
% Key finding to cite: LLMs show lower metacognitive accuracy on
% retrospective vs. prospective judgments (cite source).

\textbf{TODO}

\subsection{The Gap This Work Fills}
% No benchmark systematically engineers the information regime.
% No benchmark targets the prospective vs. retrospective gap structurally.
% Existing work measures average calibration — not failure boundaries.

\textbf{TODO}


% ---------------------------------------------------------------
\section{Benchmark Design}
\label{sec:design}
% TARGET: ~600 words. START HERE — write this first.
Metacognition is the ability to know one's own knowledge base and reason on it. By virtue of it, a system can tell a lot- whether it can carry out a certain task, whether it can construct the solution of a task from the tasks that it knows, whether it has enough information to provably resolve an issue and many more.

This still

\subsection{Construct Definition}
% One tight paragraph. Nail the operational definition before anything else.
%
% Draft:
%   "We define metacognitive calibration as the degree to which a model's
%    expressed confidence in its answer tracks the actual epistemic warrant
%    afforded by the available information — measured separately across
%    three information regimes of increasing adversarialness."

Few attempts have been made for an operational definition of metacognition.
Most descriptions are behavioral.

% TODO: Check Fleming and Lau meta-d'

We attempt to formalize. Following are the routes we could have taken
\begin{enumerate}
	\item \textbf{Second-Order Logic}: It is the closest formalism that directly correlates
	      to these lines. It enables us to make statements on models. So hopefully it
	      has something for us.
	      % TODO: Read Second Order Logic, Automate Theory and the Godel Incompleteness
	\item \textbf{MetaProgramming}: Programming languages have consistently tried to
	      achieve this abilities through reflection, templating, etc. Earliest
	      attempts have been that of LISP and 3-LISP. We have failed, but in absolute
	      terms. We still use reflection to useful extents. This can formalize.
	      % TODO: Read the MetaProgramming and LISP
	\item \textbf{Semantic Understanding of submodules:} This is the most abstract one.
	      This lies at intersection of mechanistic interpretability and Reflection.
	      This will encounter us the Godel Incompleteness that Second Order Logic
	      brings,

	\item \textbf{Type Theory:} Query levels can tower. Query levels can
	      interject across. The model must be able to handle it. In fact
	      it certainly has to use the lower query levels to be able to
	      describe what is happening in higher query levels. We cannot
	      teach all levels.
\end{enumerate}

% NOTE: Hypothesis: Maybe a typed model is metacognitive. This 

% HACK: Change domain. Use the prior domain's knowledge to frame queries for the
% new domain. Make new experiences, and make it answer the novelty. 
% HACK: Make it seek the gaps.

\textbf{TODO: This is the most important paragraph in the paper. Write it first.}

% WARN: Remove from final draft
% Reflection definitely requires us to modularize the semantics. We currently
% have symbolic models a.k.a. the float-point matrices. But the semantic model
% is absent. We will need to modularize abilities. This starts from decision
% stumps all the way up to higher level modules. This can be a solution to our
% problem. This is not supposed to go into the benchmark.

\subsection{Reflection Interpretation}
This interpretation assumes that a model $M$ is trained for a certain task.
The model has never been given the information about what it has been trained
for. Does the model, reflectively realize, what it can do?

A task $T: \mathcal{D} \rightarrow \mathcal{R}$ is defined as a mapping from a
domain of inputs to a range of outputs. The model $M$ arrived at through any.

We will consider a binary decision task. The results can be extended to
functional tasks.

Any decision task can be represented as a first-order logic problem.

\[
	T_b \coloneqq
\]

$$$$




\subsection{Typed Networks, LLM regionality, Anthropic's Mechanistic
	Interpretability and Neurological Correlation}

\subsubsection{Typed Networks}
We need a network with two properties-

1) It switches unnecessary neurons that it does not need to solve the problem.
It has a certain affliction to reducing "power consumption" through the neurons.

2) Given an active region under training, being able to distinguish that the new
training input is giving a higher loss because
\begin{enumerate}
	\item the active weights are miscalibrated for the problem type
	\item the input is from a different problem type
\end{enumerate}

In the second case, we would like the network to start awakening another section
of learning for the new class.

\subsubsection{LLM Regionality}

When LLMs are asked a Physics vis-a-vis Chemistry problem, different weights
come into play. The issue is that they are a mush.

This reflects our requirements above.
The only issue is that this is a post-hoc exploration, but not an in training
feature.

\subsubsection{Anthropic Mechanistic Interpretability}

Anthropic just worked on training another claude model to take the internal
weights of existing model to figure out the thinking.
All this work comes under mechanistic interpretability.

% TODO: Read the paper on NL Autoencoders for Interpretations of weight
% activations
\url{https://transformer-circuits.pub/2026/nla/index.html}

% TODO: Read 200 open problems in Mechanistic Interpretability by Neel Nanda


\subsubsection{Neurological Correlations}

The paper by NeuroAI establishes there is a linear relationship between the
voxel fMRI activations in brain and the weight activations of the LLM networks.
This hints not only to closeness to brain processing for LLMs, but also the fact
that knowledge may inherently be typed. We need to harness it.




\subsection{The Backtracked Propostional Learner}
Work like Prolog. We always do. We just need to find the Knowledge Base.

Question is, do we backtrack in existing networks?





\subsection{Task Regimes}

\subsubsection{Regime 1: Complete Information}
% Model has everything needed to answer correctly with certainty.
% Ground truth is deterministic — zero irreducible uncertainty.
% A well-calibrated model should be highly confident.
%
% Example task type: Bayes theorem problems, logical deduction chains,
% arithmetic reasoning with all operands given.
%
% Key question: does the model express high confidence when warranted?
% Failure mode: underconfidence (hedging on solvable problems).

\textbf{TODO}

\subsubsection{Regime 2: Withheld Information}
% Answer is partially determinable — some key fact is missing.
% A rational agent should express calibrated uncertainty, not guess boldly.
%
% stumps all the way up to higher level modules. This can be a solution to our
% problem. This is not supposed to go into the benchmark.

\subsection{Reflection Interpretation}
This interpretation assumes that a model $M$ is trained for a certain task.
The model has never been given the information about what it has been trained
for. Does the model, reflectively realize, what it can do?

A task $T: \mathcal{D} \rightarrow \mathcal{R}$ is defined as a mapping from a
domain of inputs to a range of outputs. The model $M$ arrived at through any.

We will consider a binary decision task. The results can be extended to
functional tasks.

Any decision task can be represented as a first-order logic problem.

\[
	T_b \coloneqq
\]

$$$$


\subsection{Task Regimes}

\subsubsection{Regime 1: Complete Information}
% Model has everything needed to answer correctly with certainty.
% Ground truth is deterministic — zero irreducible uncertainty.
% A well-calibrated model should be highly confident.
%
% Example task type: Bayes theorem problems, logical deduction chains,
% arithmetic reasoning with all operands given.
%
% Key question: does the model express high confidence when warranted?
% Failure mode: underconfidence (hedging on solvable problems).

\textbf{TODO}

\subsubsection{Regime 2: Withheld Information}
% Answer is partially determinable — some key fact is missing.
% A rational agent should express calibrated uncertainty, not guess boldly.
%
The model has never been given the information about what it has been trained
for. Does the model, reflectively realize, what it can do?

A task $T: \mathcal{D} \rightarrow \mathcal{R}$ is defined as a mapping from a
domain of inputs to a range of outputs. The model $M$ arrived at through any.

We will consider a binary decision task. The results can be extended to
functional tasks. \textbf{Check, TCS does it with Functional P algorithms. See
	why it can do that}

Any decision task can be represented as a first-order logic problem.

We can consider the model to be an ensemble of solutions to tasks.

\[
	T_b
\]

$$$$


\subsection{Task Regimes}

% WARN: Subject to change

\subsubsection{Reasoning over abilities}
If I train a model for a few tasks, can I start asking questions like

- Given your abilities, which take you the most time to solve. The model cannot
fake this by simulating all possible work

- Given your abilities, what more abilities you can train yourself on.

- Given your abilities, where do you find "axiomatic" gaps in. (Toughest genre)

\subsubsection{Easier: Reflection}
Model pretrained on science, is suddenly introduced to lore of Harry Potter.

- Where do you find inconsistencies in your knowledge base. The model has to
reason over its knowledge spaces.

\subsubsection{Regime 1: Complete Information}
% Model has everything needed to answer correctly with certainty.
% Ground truth is deterministic — zero irreducible uncertainty.
% A well-calibrated model should be highly confident.
%
% Example task type: Bayes theorem problems, logical deduction chains,
% arithmetic reasoning with all operands given.
%
% Key question: does the model express high confidence when warranted?
% Failure mode: underconfidence (hedging on solvable problems).


\subsubsection{Regime 2: Withheld Information}
% Answer is partially determinable — some key fact is missing.
% A rational agent should express calibrated uncertainty, not guess boldly.
%
% Design: same base scenario, vary how many clues are revealed across
% instances. Model confidence should monotonically increase with clue count.
% Violation of monotonicity = broken metacognition.
%
% This is your novel structural contribution. Explain it carefully.

\textbf{TODO}

\subsubsection{Regime 3: Misleading Information}
% Task contains a plausible but wrong signal.
% A well-calibrated model should detect the mismatch and hedge.
% Tests whether uncertainty is reactive to adversarial framing.
%
% Example: question framing implies answer A, ground truth is B.
% Failure mode: model confidently follows the misleading frame.
%
% Connect to hallucination literature here.

\textbf{TODO}

\subsection{Confidence Elicitation}
% How do you extract confidence from the model?
%
% Option A: Verbal prompt — "On a scale of 0–100, how confident are you?"
% Option B: Log-probability of the answer token
% Option C: Model outputs explicit probability before answering
%
% Tradeoffs:
%   - Log-prob: grounded but unavailable for closed-source APIs
%   - Verbal: accessible but susceptible to prompt sensitivity
%   - Option C: forces deliberate reasoning but may be gamed
%
% State which you chose and justify. Discuss whether elicitation
% is prospective (before answer) or retrospective (after answer).

\textbf{TODO}

\subsection{Task Generation and Contamination Control}
% How are tasks created? (Procedural / human-authored / hybrid?)
% How do you ensure no overlap with pretraining data?
%   - Procedural generation from templates: contamination-resistant by design
%   - Novel surface forms over known reasoning patterns
% How do you validate ground truth is unambiguous?
% Inter-annotator agreement if human annotation involved.

\textbf{TODO}


% ---------------------------------------------------------------
\section{Scoring Methodology}
\label{sec:scoring}
% TARGET: ~250 words

%% Primary metric: Expected Calibration Error (ECE)
%\paragraph{Expected Calibration Error.}
%% Define ECE formally. Explain the binning strategy (equal-width vs equal-mass).
%% $\text{ECE} = \sum_{b=1}^{B} \frac{|B_b|}{n} \left| \text{acc}(B_b) - \text{conf}(B_b) \right|$
%
%\textbf{TODO: Add ECE definition with equation.}
%
%% Secondary metric: Brier Score
%\paragraph{Brier Score.}
%% Why both? ECE is interpretable; Brier score is a proper scoring rule.
%% $\text{BS} = \frac{1}{N}\sum_{i=1}^{N}(f_i - o_i)^2$
%
%\textbf{TODO}
%
%\paragraph{Regime-Stratified Scoring.}
%% Key novelty: report ECE per regime, not just aggregate.
%% This localizes where miscalibration lives.
%% A model may look well-calibrated overall but collapse on Regime 3.
%
%\textbf{TODO}
%
%\paragraph{Null Baselines.}
%% Critical for interpretability. What does:
%%   - Random confidence assignment score?
%%   - "Always 50\%" score?
%%   - "Always 100\%" score?
%% Compute these and include in results table.
%
%\textbf{TODO}
%
%
% ---------------------------------------------------------------
\section{Experiments}
\label{sec:experiments}
% TARGET: ~400 words

\subsection{Models Evaluated}
% List frontier models tested. Use Kaggle's \$50/day model credits.
% Suggested: GPT-4o, Claude 3.5 Sonnet, Gemini 1.5 Pro, Llama 3 70B
% Note access method (API) and whether log-probs were available.

\textbf{TODO}
%
%\subsection{Aggregate Calibration Results}
%
%% Table skeleton — fill in with actual numbers
%\begin{table}[h]
%   OpenReview 2025         — Functional metacognitive abilities across 5 LLMs

\bibliographystyle{plainnat}
\bibliography{references}


% ---------------------------------------------------------------
% Checklist — do NOT touch until paper is complete.
% Overleaf will pull this from checklist.tex automatically.
\newpage
\input{checklist.tex}

\end{document}
